% !TEX root = ../../main/main.tex
\hypearticle
{A IA irá reduzir a ciência e a educação para um algoritmo ou fará todos nós acolhermos nossa humanidade?}
{Nathália Pereira}
{

\subtitle{O impacto da IA}
    \hspace{1cm}Sibel Erduran, da Universidade de Oxford, afirma em seu artigo que a Inteligência Artificial (IA) apresenta um crescimento exponencial nos últimos anos,
de tal forma a adentrar gradualmente no ambiente científico e metodológico. Essa transformação vem a culminar em uma mudança significativa nos métodos de
ensino e aprendizagem, de forma a trabalhar no \textbf{limiar entre o auxílio e a redução do que é humano.}

\subtitle{Contextualização}
    \hspace{1cm}Ao mesmo tempo em que é apresentado uma \textbf{atuação benéfica da IA para a ciência}, ocasionando
desde descobertas até o desenvolvimento de técnicas para pesquisas, também é relatado o quanto a intervenção da IA no método científico
tem \textbf{alterado as convenções da própria ciência}, de modo a alterar meios de coleta de dados, organização e análise dos mesmos, por exemplo.

    \hspace{1cm}Apesar de haver um baixo índice de artigos que identificam a IA como influenciadora nas etapas de conhecimento e produções científicas, percebeu-se
que, com o lançamento do ChatGPT, houve uma comoção entre os cientistas. Por consequência, cada vez mais atenção foi direcionada às \textbf{dimensões
sociais}. Assim, é destacado a ética, a política e a governança institucional, como no caso do cuidado de um médico com seu paciente ao contrário da menção de sintomas ao ChatGPT.

\subtitle{Análise crítica}
    \hspace{1cm}Desse modo, é possível relacionar ao estudo que, considerando o comportamento humano imediatista, à luz da Modernidade Líquida de Bauman, o simples uso da IA
como um meio para agilizar o método científico torna o conhecimento e esforço humano altamente fluidos e mutáveis, uma vez que a maturação do aprendizado é descartada.

    \highlight{Desconhecer os procedimentos é desconhecer o produto, ou seja, uma alienação induzida e normalizada.}

\subtitle{Outra visão}
    \hspace{1cm}Por outro lado, por mais que a IA tenha recebido várias críticas, pesquisadores da educação em ciência têm dedicado seu tempo para entender como
 as plataformas de IA ainda podem ser \textbf{produtivas} para aulas de ciência, como na situação da IA generativa poder ocupar-se do trabalho manual de resumos.
Da mesma forma, percebeu-se que os chatbots de IA impulsionam a confiança do usuário, e por meio da reafirmação, diminuem seus índices de ansiedade e insegurança.
 Essa observação revela uma \textbf{implicação ética da IA na educação}, já que estimula uma maior vontade de estudo.

\subtitle{Ameaça à própria identidade humana}
    \begin{infobox} O texto aponta que a IA está provocando a indagação do que é “realmente humano”. \end{infobox}

    \hspace{1cm}Logo, questionamentos surgem: como a ciência precisa ser feita, como a educação precisa ser reconceituada e como as pesquisas necessitam ser construídas?

    \hspace{1cm}Nesse sentido, o aprendizado crescente da IA tem se assemelhado à capacidade lógica humana, de escrita, de interpretação e de criatividade. \textbf{Assim, a IA
permite uma tênue linha entre o que distingue o ser humano do resto de outros seres vivos: a capacidade cognitiva.}

\subtitle{Algoritmo x Humanidade}
    \highlight{Consequentemente, a habilidade cognitiva não pode mais ser a principal prova de ser um humano.}

    \hspace{1cm}O limiar criado entre o ser humano e a máquina gera uma busca por formas alternativas de visualizar o que é realmente e unicamente humano.

    \hspace{1cm}Assim, surge a classificação por \textbf{inteligência funcional e inteligência consciente, respectivamente, dos modelos de IA e do ser humano.}

    \hspace{1cm}Nessa visão, o autor aponta que uma máquina não apresenta a empolgação, o momento de ‘eureka’, após ter uma brilhante descoberta tal como ocorre com os
cientistas, e também não é capaz de produzir um ensino tão profundo e conexo quanto o de um professor ao seu aluno.

    \hspace{1cm}Em sequência, é possibilitado a indução de que a
IA funciona tal como um algoritmo treinado para receber uma entrada e produzir uma saída. Já os humanos são animais sociais, eles necessitam de um \textbf{toque humano
em sua realidade e de experiência: a humanidade, a cultura.}

\subtitle{Como prosseguir?}
    \hspace{1cm}No final, o trabalho atribui a responsabilidade de mediar a relação da IA com o ensino e o método científico aos futuros pesquisadores
da educação em ciências, de modo que ainda não os transforme em simples algoritmos, mas que utilize a IA como uma forma de reforçar
o que podemos fazer com nossa experiência e valores humanos dentro desses campos.

    Por fim, estendo o questionamento:
    \begin{infobox} Quando deixaremos de tratar a IA como um risco e passaremos à enxergar como oportunidade? \end{infobox}


\originalsource{https://www.scielo.br/j/ciedu/a/jJRG5nJVpKLRzX5hRxFM9ns/?lang=en}

}